% ============================================================
%  Statische Geometrie und SI-Projektion
%  Warum natuerliche Einheiten keine fraktalen Korrekturen brauchen
%  Kapitelinhalt DE
%  Quellenbelegt: Dok. 011, 012, 013, 014, 015, 041, 044, 116, 133, 193, 258
% ============================================================

\title{Dok.~261 — Statische Geometrie und SI-Projektion:\\
Warum natürliche Einheiten keine fraktalen Korrekturen brauchen}
\author{Johann Pascher}
\date{Mai 2026}
\maketitle

\begin{center}
\textit{Kontext: Dok.~011 (Feinstruktur), Dok.~013/014 (SI \& natürliche Einheiten),
Dok.~116 (Koide), Dok.~133 (fraktale Korrektur)} \\[2pt]
\textit{FFGFT-Archiv: DOI 10.5281/zenodo.20474821}
\end{center}

\tableofcontents
\newpage

%=============================================================
\section{Einordnung}
%=============================================================

Die SI-Reform von 2019 hat vier Konstanten ($c$, $\hbar$, $e$, $k_B$) auf exakte
Werte fixiert und damit die meisten dimensionsbehafteten Größen zu reinen
Einheitendefinitionen gemacht. Was als \emph{echter} empirischer Input übrig
bleibt, sind die dimensionslosen Konstanten -- allen voran die
Feinstrukturkonstante $\alpha$, deren Wert seit der Reform weiterhin gemessen
wird ($\mu_0$, $\varepsilon_0$ und $\alpha$ blieben Messgrößen). FFGFT leitet
auch diese dimensionslosen Inputs aus dem einzigen Parameter
$\xipar = \tfrac{4}{3}\times 10^{-4}$ ab.

Dieses Dokument stellt die Umrechnungssystematik kompakt zusammen
(Tabellen~\ref{tab:si-fixiert}--\ref{tab:xi-inputs}) und begründet anschließend
einen zentralen Punkt: \emph{solange man in natürlichen Einheiten und reinen
Verhältnissen bleibt, ist die Beschreibung statisch und relativ und erfordert
keine fraktalen Korrekturen.} Die fraktalen Korrekturen treten erst beim Übergang
in absolute SI-Werte auf.

%=============================================================
\section{Kompaktifizierte Zeit: statisch und vollständig --- der Pfeil macht dynamisch}
\label{sec:rahmen}
%=============================================================

Dieser Abschnitt stellt den Rahmen über alle folgenden: Ob eine Größe statisch
oder dynamisch erscheint, hängt allein vom \emph{Status der Zeitrichtung} ab --
nicht von der Zahl der Dimensionen.

Die FFGFT-Vollgeometrie ist der 4-Torus $T^4 = T^3 \times S^1$, in dem die Zeit
eine \emph{kompakte} Richtung ist -- die Zeit-Wicklung $S^1$ (Dok.~230). In dieser
vollen Sicht ist die Zeit eine geschlossene Schleife, gleichberechtigt mit den
drei räumlichen Wicklungen; es gibt kein fließendes ``Jetzt'' und keinen Pfeil.
Die Beschreibung ist rein relational: Windungszahlen und dimensionslose
Verhältnisse. Das ist die \emph{Schicht~1} aus Dok.~241 -- die Ebene, auf der
FFGFT ihre fundamentalen Aussagen macht.

Auf dieser Schicht wirkt alles statisch, und doch ist nichts verloren. Die
fraktale Rekursion der Theorie ist hier nicht abwesend, sondern \emph{eingefroren}:
Ein Verhältnis wie $m_\mu/m_e \approx 206{,}768$ ist das Endprodukt einer
Rekursion, die in dieser Zahl zur Ruhe gekommen ist -- so wie der goldene Schnitt
$\varphi$ nicht \emph{neben} der Fibonacci-Folge steht, sondern ihr Grenzwert
\emph{ist} (Dok.~241). Was auf Schicht~1 fehlt, ist nicht die Fraktalität, sondern
ihre sichtbare Bewegung.

\begin{prinzip}
\textbf{Statisch ist nicht unvollständig.} Im kompakten $T^4$ (Zeit als Wicklung)
trägt die Beschreibung allen invarianten Inhalt -- als eingefrorene Fixpunkte
fraktaler Rekursionen. Der Übergang in die Dynamik fügt keinen Inhalt hinzu; er
macht nur die Rekursion sichtbar in Bewegung.
\end{prinzip}

Dynamik entsteht durch \emph{Dekompaktifizierung}: Wird die Zeit-Wicklung $S^1$ zum
offenen Pfeil $t$ aufgelöst -- geometrisch der Schritt $T^4 \to T^3$, der ``die
Zeit-Wicklung verliert'' (Dok.~230) --, so bleiben drei Raumrichtungen plus eine
gerichtete Zeit: die beobachtete 3+1-Welt. Was diese eine Richtung als Pfeil
heraushebt, ist die fundamentale $\xipar$-Asymmetrie, die ``alle
Entwicklungsprozesse ermöglicht'' (Dok.~078). Erst jetzt läuft die Rekursion
sichtbar: Prozesse entfalten sich, Phase akkumuliert, und die
\emph{Einbettungskorrektur} wird explizit -- der ``Preis'', von dem Dok.~185
spricht.

Damit fügt sich alles Folgende zu einem Bild. Die Brücke von der statischen
Schicht~1 in die dynamische 3+1-Welt -- die \emph{Schicht~2}, die SI-Übersetzung
(Dok.~241) -- wird genau von $c$, $\hbar$, $\varepsilon_0$ getragen. Deshalb werden
diese Konstanten erst hier real (Abschnitt~\ref{sec:dyn-konstanten}), und deshalb
sind auch $G$ und $\alpha$ keine statischen Grundgrößen, sondern dynamische
Kopplungen (Abschnitt~\ref{sec:G-alpha}). Die einzige statische Grundgröße bleibt
$\xipar$; alles andere ist entweder ein eingefrorener $\xipar$-Fixpunkt
(Schicht~1) oder dessen sichtbar bewegte Projektion (Schicht~2).

Eine Präzisierung, damit ``vollständig'' nicht missverstanden wird: Es bedeutet
\emph{enthält allen invarianten Inhalt}, nicht \emph{die Dynamik sei unwirklich}.
Für den in den Pfeil eingebetteten Beobachter ist der Zeitfluss operativ real. Die
saubere Aussage lautet: Das kompakte $T^4$-Bild ist inhaltlich vollständig; die
Dynamik ist seine reale, aber inhaltsneutrale Innenperspektive.

Schließlich schließt sich der Bogen zur Retrokausalität. Auf dem $T^4$ mit intakter
Zeit-Wicklung ist eine geschlossene zeitartige Kurve nur eine Windung der
Zeitschleife, und Verschränkung erscheint als \emph{topologische Verbindung}
(Dok.~230) -- statische Geometrie. Erst nach der Dekompaktifizierung in die
3+1-Pfeilzeit sieht dieselbe Geometrie wie Rückwärtskausalität oder Nichtlokalität
aus. Die scheinbare Exotik des Retrokausalen ist damit ein Artefakt der
dynamischen Innensicht, nicht der Geometrie.

%=============================================================
\section{Die durch die SI-Reform fixierten Konstanten}
%=============================================================

\begin{table}[htbp]
  \centering
  \caption{Die durch die SI-Reform 2019 fixierten Konstanten (die ``Maschinerie'').
           Sie tragen keine Physik, nur Maßstab, und werden in natürlichen
           Einheiten zu~$1$.}
  \label{tab:si-fixiert}
  \adjustbox{max width=\textwidth}{%
\begin{tabular}{llll}
    \toprule
    Konstante & Wert (exakt seit 2019) & definiert & nat.\ Einheiten \\
    \midrule
    $c$     & $299\,792\,458$ m/s                          & Meter   & $=1$ \\
    $\hbar$ & $1{,}054\,571\,817\ldots\times10^{-34}$ J\,s & kg      & $=1$ \\
    $e$     & $1{,}602\,176\,634\times10^{-19}$ C           & Ampere  & Ladungseinheit \\
    $k_B$   & $1{,}380\,649\times10^{-23}$ J/K              & Kelvin  & $=1$ \\
    \bottomrule
  \end{tabular}%
}
\end{table}

%=============================================================
\section{Größenumrechnung in natürlichen Einheiten}
%=============================================================

\begin{table}[htbp]
  \centering
  \caption{Größenumrechnung in natürlichen Einheiten ($\hbar=c=k_B=1$):
           alle dimensionsbehafteten Größen werden zu Potenzen der Energie.}
  \label{tab:nat-umrechnung}
  {\small
  \adjustbox{max width=\textwidth}{%
\begin{tabular}{lccl}
    \toprule
    Größe & SI-Dimension & $[\hbar=c=k_B=1]$ & SI-Rückrechnung \\
    \midrule
    Energie       & J        & $E$           & Basis \\
    Masse         & kg       & $E$           & $\times\,S_{\mathrm{T0}}$ \ (über $c^{-2}$) \\
    Impuls        & kg\,m/s  & $E$           & $\times\,c^{-1}$ \\
    Länge         & m        & $E^{-1}$      & $\ell[\mathrm{fm}]=197{,}327/E[\mathrm{MeV}]$ \\
    Zeit          & s        & $E^{-1}$      & $t[\mathrm{s}]=6{,}582\times10^{-22}/E[\mathrm{MeV}]$ \\
    Temperatur    & K        & $E$           & $T[\mathrm{K}]=E[\mathrm{eV}]/(8{,}617\times10^{-5})$ \\
    Ladung (Gauß) & C        & dimensionslos & $e=\sqrt{\alpha}\approx0{,}0854$ \\
    \bottomrule
  \end{tabular}%
}}

  \medskip
  \footnotesize
  Die SI-Rückrechnung läuft über die \emph{eine} Massenskala
  $S_{\mathrm{T0}} = 1\,\mathrm{MeV}/c^2 = 1{,}782662\times10^{-30}$~kg (Dok.~013/014);
  der Längenfaktor ist $\hbar c/S_{\mathrm{T0}} = 1{,}973\times10^{-13}$~m, also die
  $197$\,MeV\,fm aus der $\hbar c=1$-Spalte in T0-Sprache.\par\smallskip
  In T0-Einheiten wird die Ladungseinheit so gewählt
  ($e=\sqrt{4\pi\varepsilon_0\hbar c}$), dass der \emph{geschriebene} Wert
  $\alpha=1$ ist (Dok.~044) -- eine Konvention, kein physikalischer Eingriff.
  Die physikalische Invariante $\alpha\approx1/137$ wird aus $\xipar$ rekonstruiert
  (Tab.~\ref{tab:xi-inputs}).
\end{table}

%=============================================================
\section{Die dimensionslosen Inputs und ihre Herleitung aus $\xipar$}
%=============================================================

\begin{table}[htbp]
  \centering
  \caption{Die überlebenden dimensionslosen Inputs der Physik und ihre
           Herleitung aus dem einzigen Parameter $\xipar=\tfrac{4}{3}\times10^{-4}$.
           Sämtliche Relationen quellenbelegt.}
  \label{tab:xi-inputs}
  {\footnotesize\setlength{\tabcolsep}{4pt}
  \adjustbox{max width=\textwidth}{%
\begin{tabular}{lccl}
    \toprule
    Invariante & Standardphysik & FFGFT-Relation (aus $\xipar$) & Quelle \\
    \midrule
    $\xipar$           & ---            & $=\tfrac{4}{3}\times10^{-4}=4/30000$ & Dok.~013 \\
    $\alpha$ (SI)      & $1/137{,}036$  & $\xipar\,(\Ezero/1\,\mathrm{MeV})^2$\,$^{\dagger}$ & Dok.~011 \\
    Leptonmassen       & gemessen       & $m=r\,\xipar^{\,p}\,v$ & Dok.~116, 006 \\
    \quad Exponenten   &                & $(p_e,p_\mu,p_\tau)=(\tfrac{3}{2},\,1,\,\tfrac{2}{3})$ & Dok.~116 \\
    \quad Verhältnisse &                & $(r_e,r_\mu,r_\tau)=(\tfrac{4}{3},\,\tfrac{16}{5},\,\tfrac{25}{9})$ & Dok.~116 \\
    Koide $Q$          & $0{,}66666$    & $=\tfrac{2}{3}$ \ (Folge der $\xipar$-Exponenten) & Dok.~116, 258 \\
    $\sin^2\theta_W$   & $0{,}2312$     & $\dfrac{1-\sqrt{1-4\alpha_W}}{4}$,\ \ $\alpha_W=\sqrt{\xipar}$ & Dok.~041 \\[0.6em]
    $\theta_{13}$      & $8{,}57^\circ$ & $\arcsin\!\big(\xipar^{1/3}\big)$ & Dok.~041 \\
    $G$                & gemessen       & $\dfrac{\xipar^2}{4m_e}\,\Kfrak$,\ \ $\Kfrak=1-100\xipar\approx0{,}9867$ & Dok.~012, 133 \\
    \bottomrule
  \end{tabular}%
}}

  \medskip
  \footnotesize
  $^{\dagger}$\,$\alpha=\xipar\,(\Ezero/1\,\mathrm{MeV})^2$ ist die \emph{Leitordnung},
  mit $\Ezero=\sqrt{m_e\,m_\mu}=7{,}398$\,MeV.
  Die CODATA-Präzision wird über das fraktale Korrekturprodukt
  $\displaystyle\prod_{n=1}^{137}\!\big(1+\delta_n\,\xipar\,(4/3)^{\,n-1}\big)$
  erreicht. Diese Korrekturen sind \textbf{keine freien Annahmen}:
  $\Kfrak=1-100\xipar$ ist über den Skalenfluss von der Planck- zur
  elektroschwachen Skala eindeutig fixiert ($D_f^{\,\mathrm{Raum}}=3-\xipar$,
  Dok.~133, ``$D_f$ ist eindeutig bestimmt -- nicht frei wählbar''), und die
  rationalen Leiter-Verhältnisse $r_i$ sind $\alpha$-geometrische Invarianten
  (Dok.~258), strukturell unzugänglich für nachträgliche Anpassung
  (Nichtabschluss-Theorem, Dok.~193).
\end{table}

%=============================================================
\section{Statische Geometrie: warum natürliche Einheiten korrekturfrei sind}
%=============================================================

Solange die Physik in natürlichen Einheiten und reinen Verhältnissen ausgedrückt
wird, ist ihre Beschreibung \emph{statisch} und \emph{relational}: Es treten nur
dimensionslose Verhältnisse auf, keine absoluten Anker und keine
Skalenabhängigkeit. Auf dieser Ebene sind die FFGFT-Größen exakte geometrische
Invarianten von $\xipar$ -- und es ist keine fraktale Korrektur erforderlich.

\begin{prinzip}
\textbf{Prinzip.} In natürlichen Einheiten ist FFGFT rein relational. Die
dimensionslosen Strukturrelationen -- Massenverhältnisse, Mischungswinkel, die
Koide-Größe -- sind statische geometrische Invarianten von $\xipar$. Sie
enthalten keine fraktalen Korrekturen.
\end{prinzip}

Das deutlichste Beispiel ist die Koide-Größe: $Q=\tfrac{2}{3}$ folgt exakt aus den
Exponenten $(p_e,p_\mu,p_\tau)=(\tfrac{3}{2},1,\tfrac{2}{3})$ und den
Verhältnissen $(r_e,r_\mu,r_\tau)=(\tfrac{4}{3},\tfrac{16}{5},\tfrac{25}{9})$ --
ohne zusätzliche Parameter und ohne fraktale Korrektur (Dok.~116). Ebenso sind die
Leiter-Verhältnisse $r_i$, die Exponenten $p_i$, die Weinberg-Struktur und
$\theta_{13}$ exakte Funktionen von $\xipar$ allein. Keine dieser Relationen
benötigt einen absoluten Maßstab, ein $\Ezero$ in MeV oder einen
Umrechnungsfaktor.

Wo entstehen dann die fraktalen Korrekturen? Erst an der \emph{SI-Projektion} --
dort, wo eine absolute Skala festgelegt wird. $\Kfrak=1-100\xipar$ ist nicht eine
Korrektur an der Geometrie, sondern die kumulative Wirkung des Skalenflusses von
der Planck- zur elektroschwachen Skala (Dok.~133): eine Eigenschaft der Abbildung
der Geometrie auf ein absolutes Energie-Lineal. Genauso erscheint das fraktale
Korrekturprodukt $\prod_n(1+\delta_n\,\xipar\,(4/3)^{n-1})$ erst, wenn die führende
Beziehung $\alpha=\xipar(\Ezero/1\,\mathrm{MeV})^2$ auf CODATA-Präzision in
absoluten SI-Werten gebracht wird.

Wichtig zur Abgrenzung: Diese Projektionsfaktoren sind dennoch \emph{keine} freien
Parameter. $\Kfrak$ ist über den Skalenfluss eindeutig fixiert (Dok.~133), die
rationalen Verhältnisse sind $\alpha$-geometrische Invarianten (Dok.~258), und das
Nichtabschluss-Theorem (Dok.~193) sperrt jede nachträgliche Anpassung. Sie sind
abgeleitete Projektionsfaktoren, nicht angenommene Fits.

\begin{schlussfolgerung}
\textbf{Schlussfolgerung.} Die fraktalen Korrekturen sind der Preis der
Absolut-Projektion, nicht ein Defekt der Geometrie. Solange man in natürlichen
Einheiten und Verhältnissen bleibt, ist die Beschreibung statisch, relativ und
korrekturfrei; die Korrekturen treten ausschließlich beim Übergang in absolute
SI-Werte auf -- und sind selbst dort eindeutig aus $\xipar$ und der
Skalenstruktur bestimmt.
\end{schlussfolgerung}

%=============================================================
\section{Dynamische Realität der Naturkonstanten: $c$, $\hbar$, $\varepsilon_0$ im Zeitfluss}
\label{sec:dyn-konstanten}
%=============================================================

Der vorige Abschnitt zeigt, dass die statische $\xipar$-Geometrie korrekturfrei
ist. Es bleibt die Frage, \emph{wann} die Konstanten $c$, $\hbar$ und
$\varepsilon_0$ überhaupt physikalische Realität annehmen. Die Antwort folgt aus
derselben Trennung -- nur von der anderen Seite betrachtet: Diese Konstanten sind
keine Eigenschaften der ruhenden Geometrie, sondern die Umrechnungsraten, die sie
annimmt, sobald sie in einen realen, zeitlich ausgedehnten Prozess eingebettet
wird.

FFGFT trennt dazu zwei Zeiten (Dok.~078): die intrinsische, charakteristische
Skala $\tilde T$ der Dualität $\tilde T\cdot m=1$, die diskrete
Systemeigenschaften reguliert, und den übergeordneten, kontinuierlichen Zeitfluss
(die ``Pfeilzeit'' $t$), in dem die Systeme existieren und Prozesse ablaufen. Die
Dualität betrifft die intrinsische Skala, nicht den Zeitfluss selbst. Erst im
Zeitfluss -- wenn eine zeitlose Identität physikalisch \emph{hergestellt} wird --
entsteht das, was Dok.~185 den \emph{Einbettungspreis} nennt: Die Realisierung
einer statischen mathematischen Tatsache als Prozess kostet reale Zeit. $c$,
$\hbar$ und $\varepsilon_0$ sind die Wechselkurse dieses Einbettungspreises.

\begin{erkenntnis}
\textbf{Einbettungspreis.} Eine zeitlose $\xipar$-Relation ist eine statische
Tatsache. Ihre physikalische Realisierung als Prozess in der Pfeilzeit kostet
Zeit -- und $c$, $\hbar$, $\varepsilon_0$ sind genau die Raten, die diesen Preis
quantifizieren (Dok.~185).
\end{erkenntnis}

Im Einzelnen:

\textbf{$c$} ist kein fundamentales Naturgesetz, sondern ein dynamisches
Verhältnis $L/T$ (Dok.~134). Statisch sind Energie und Masse dieselbe Zahl
($E=m$, $c=1$, $\text{Länge}\equiv\text{Zeit}$); $c$ erscheint gar nicht. Sobald
sich etwas ausbreitet, wird $c$ die Durchlaufrate der $T^4$-Geometrie -- die volle
Form $E=mc^2$, die Dispersion $E^2=(pc)^2+(mc^2)^2$, der Lichtkegel. $c$ ist real
als die Rate, mit der die Geometrie zeitlich abgefahren wird.

\textbf{$\hbar$} ist der Wechselkurs zwischen geometrischer Windung und
akkumulierter realer Zeit. Statisch ist Phase eine dimensionslose Windungszahl auf
dem Torus ($\hbar=1$). Im Zeitfluss akkumuliert die Phase $e^{-iEt/\hbar}$;
$\hbar$ ist die Rate, mit der die geometrische Windung als Oszillation,
Interferenz und Zerfall sichtbar wird. In der Feldfassung $T_{\text{Feld}}\cdot
E_{\text{Feld}}=1$ (Dok.~020) ist dies der Übergang von der intrinsischen Skala
zur beobachtbaren Dynamik.

\textbf{$\varepsilon_0$} ist die Antwort des Vakuums auf zeitveränderliche Felder.
Statisch steckt sie in der Ladungseinheiten-Konvention $\alpha=1$. Dynamisch wird
sie real über $c=1/\sqrt{\varepsilon_0\mu_0}$: $\varepsilon_0$ und $\mu_0$
resonieren dual (Vakuum als resonantes Medium, Dok.~043/044), und ihr Produkt
\emph{ist} die Lichtausbreitung in der Zeit. $\varepsilon_0$ ist real als
Impedanz des Vakuums gegen den zeitlichen Feldfluss.

Die Konsequenz für die Teilchenphysik ist eine scharfe Trennlinie:

\begin{prinzip}
\textbf{Demarkation.} \emph{Dimensionslose, statische} Größen
(Massenverhältnisse, Mischungswinkel, $\sin^2\theta_W$, Koide $Q=2/3$) leben
vollständig in der $\xipar$-Geometrie -- ohne $c$, $\hbar$, $\varepsilon_0$ und
ohne fraktale Korrektur. \emph{Absolute, zeitausgedehnte} Größen (Lebensdauern,
Zerfallsraten, Wirkungsquerschnitte, Oszillationsfrequenzen, Propagatoren)
existieren erst im Zeitfluss, erfordern $c$, $\hbar$, $\varepsilon_0$ als real und
erben damit die SI-Projektion ($S_{\mathrm{T0}}$, $\Kfrak$).
\end{prinzip}

Das ist konkret testbar: Ein Lebensdauer-\emph{Verhältnis} (zwei Prozesse,
$\hbar$ kürzt sich) ist statisch-sauber aus $\xipar$ ableitbar (vgl.~Dok.~160);
eine \emph{absolute} Lebensdauer trägt die dynamischen Konstanten und die
Projektionsfaktoren. Die Demarkation sagt also voraus, welche Größen FFGFT scharf
und korrekturfrei liefert und welche die Projektion durchlaufen müssen.

Eine ontologische Präzisierung zum Abschluss: ``real'' heißt hier \emph{operativ
real im dynamischen Regime}, nicht \emph{ontologisch fundamental}. Die Geometrie
($\xipar$, $T^4$, $\tilde T\cdot m=1$) ist primär; $c$, $\hbar$, $\varepsilon_0$
sind die realen Umrechnungsraten, die das zeitlose Verhältnisgefüge annimmt,
sobald es prozesshaft wird. Sie sind die SI-Schatten der Einbettung der statischen
Geometrie in den Zeitfluss.

%=============================================================
\section{Konsequenz: $G$ und $\alpha$ sind keine statischen Grundkonstanten}
\label{sec:G-alpha}
%=============================================================

Genau dieselbe Demarkation trifft die beiden Größen, die traditionell als
fundamentale Konstanten gelten: die Gravitationskonstante $G$ und die
Feinstrukturkonstante $\alpha$. Beide sind \emph{Kopplungskonstanten} -- sie
quantifizieren die Stärke einer Wechselwirkung. Und eine Kopplungsstärke ist der
klarste Fall einer dynamischen Größe: sie beschreibt einen Prozess und sie
\emph{läuft} mit der Skala.

\textbf{$\alpha$.} Strukturell ist $\alpha$ kein freier Input, sondern aus
$\xipar$ abgeleitet: $\alpha=\xipar\,(\Ezero/1\,\mathrm{MeV})^2$ (Dok.~011,
vgl.~Tab.~\ref{tab:xi-inputs}). Operativ aber ist
$\alpha=e^2/(4\pi\varepsilon_0\hbar c)$ -- gebaut aus genau den
Einbettungs-Konstanten $\varepsilon_0$, $\hbar$, $c$ -- und sie \emph{läuft} mit
der Energie ($\alpha(Q^2)$). Dieses Laufen ist nichts anderes als der
Skalenfluss, der auch das fraktale Korrekturprodukt
$\prod_n(1+\delta_n\,\xipar\,(4/3)^{n-1})$ erzeugt. $\alpha$ ist also weder
fundamental (sie folgt aus $\xipar$) noch statisch (sie lebt im
Skalen-/Zeitfluss).

\textbf{$G$.} In natürlichen Einheiten ist $G=1$ (Dok.~012). Der SI-Wert folgt aus
$\xipar$: $G=\dfrac{\xipar^2}{4m_e}\,\Kfrak$ -- und enthält damit den
Projektionsfaktor $\Kfrak=1-100\xipar$, der selbst ein Skalenfluss-Effekt ist
(Dok.~133). Die äquivalente Planck-Form $G=\ell_{\mathrm P}^2 c^3/\hbar$
(Dok.~012/127) zeigt es noch direkter: $G$ ist durch $c^3/\hbar$ ausgedrückt, also
durch die dynamischen Konstanten selbst. $G$ ist damit doppelt projiziert -- über
$\Kfrak$ und über $c$, $\hbar$ -- und misst die Stärke der gravitativen
Wechselwirkung, eines Prozesses. Auch $G$ ist weder fundamental noch statisch.

\begin{schlussfolgerung}
\textbf{Schlussfolgerung.} Die einzige statische Grundgröße ist $\xipar$. $G$ und
$\alpha$ sind keine fundamentalen, statischen Konstanten, sondern \emph{dynamische
Kopplungen}: aus $\xipar$ abgeleitet, mit den Einbettungs-Konstanten $c$, $\hbar$,
$\varepsilon_0$ angezogen und durch Skalenfluss-Faktoren ($\Kfrak$, das fraktale
Produkt) gekleidet. Sie werden -- wie $c$, $\hbar$, $\varepsilon_0$ -- erst im
Zeitfluss real. Die scheinbare Fundamentalität von $G$ und $\alpha$ ist eine Folge
der Messkonvention, kein Fundament.
\end{schlussfolgerung}

%=============================================================
\section{Die Demarkation aus Lagrange-Sicht}
%=============================================================

Die statisch/dynamische Trennung dieses Dokuments ist keine Lesart der
Tabellen~\ref{tab:si-fixiert}--\ref{tab:xi-inputs}, sondern die Architektur der
Lagrange-Erweiterung selbst. Die Tabellen fassen nur zusammen, was der Lagrangian
bereits trägt.

\textbf{Zwei Dichten als Notwendigkeit.} Dok.~095 zeigt, dass FFGFT zwei
komplementäre Lagrange-Dichten verlangt -- nicht aus Stil, sondern aus der
Skalenhierarchie. Die vereinfachte
\[
  \mathcal{L}_{\mathrm{T0}} = \varepsilon\,(\partial\,\delta E)^2,
  \qquad \varepsilon = \xipar\,E^2
\]
trägt reine $\xipar$-Struktur (Schicht~1); die erweiterte Standardmodell-Dichte mit
T0-Korrekturen trägt die Dynamik (Schicht~2). Beide sind, wie Dok.~095 betont,
``mathematische Beschreibungen ... nicht die Natur selbst''. Die Demarkation ist
damit ein Korollar der zwei Dichten, kein interpretativer Überbau.

\textbf{Der Schnitt verläuft Term für Term.} Die vollständige Zerlegung (Dok.~067)
trennt Ableitungs- von Massentermen:
\[
  \mathcal{L}_{\text{Zeit}} = \sqrt{-g}\Big[\tfrac{1}{2}\,g^{\mu\nu}\partial_\mu T\,\partial_\nu T - V(T)\Big],
  \qquad
  \mathcal{L}_{\text{Eich}} = \sqrt{-g}\big(-\tfrac{1}{4}F_{\mu\nu}F^{\mu\nu}\big),
\]
\[
  \mathcal{L}_\psi = \sqrt{-g}\,\Omega^4(T)\big(i\bar\psi\gamma^\mu\partial_\mu\psi - m\,\bar\psi\psi\big).
\]
Die \emph{Ableitungsterme} ($\partial_\mu$, $\sqrt{-g}$ mit der Metrik,
$-\tfrac14 F^2$) tragen die Dynamik -- und in dimensionsbehafteter Form $c$ (über
die Metrik) und $\hbar$ (über die kinetische Normierung). Die
\emph{Massen-/Potentialterme} ($V(T)$, $m\bar\psi\psi$) sind statisch. Der
Euler-Lagrange-Schritt, der hieraus Bewegungsgleichungen macht, \emph{ist} die
Dekompaktifizierung in die Pfeilzeit aus Abschnitt~\ref{sec:rahmen}.

\textbf{$\alpha$ und $G$ im Lagrangian.} Die Eichkopplung sitzt in der kovarianten
Ableitung $D_\mu = \partial_\mu + ieA_\mu$ (Dok.~019) -- $\alpha$ ist also
strukturell ein Ableitungs-(Dynamik-)Term, kein statischer Parameter. Und die
Gravitation ist nicht eingesetzt, sondern \emph{emergent}: Dok.~049 zeigt, dass sie
aus $\mathcal{L} = \varepsilon\,(\partial\,\delta m)^2$ automatisch entsteht, und
Dok.~180 leitet $G$ in Schritt~4 einer Selbstkonsistenz-Kette direkt aus $\xipar$
ab -- $\xipar$, $L_0$ und $G$ werden nicht postuliert, sie folgen aus der
Selbstkonsistenz des Lagrangians.

\begin{result}
\textbf{Feldtheoretischer Ursprung.} Die Demarkation der
Abschnitte~\ref{sec:rahmen}, \ref{sec:dyn-konstanten} und \ref{sec:G-alpha} ist
eine Konsequenz der Lagrange-Erweiterung: die zwei Dichten (Dok.~095) sind
Schicht~1 und Schicht~2; der Schnitt zwischen Ableitungs- und Massentermen
(Dok.~067) ist der Schnitt statisch/dynamisch; $\alpha$ lebt in der kovarianten
Ableitung (Dok.~019); $G$ emergiert aus $\xipar$ (Dok.~049/180). Die Tabellen sind
die SI-Zusammenfassung, nicht die Beweisquelle.
\end{result}
